\section{Limitations and Open Work}
\label{sec:limitations}

We state plainly what the current evidence does and does not support.

\begin{itemize}
  \item \textbf{Real-data acceptance remains limited.} A specified-source UCI red-wine pilot reached a local SA
        solution both through saved-model recovery and a fresh run, with independent case-specific checks.
        Recovered wholesale-customer and real-network cases also delivered independently checked solutions.
        A ten-case real-data acceptance catalog has been prepared but not fully executed. The new autonomous
        batch suite uses Taipei chargers, YouBike routing and AED grouping (\cref{sec:autonomousprotocol});
        it has development smoke tests but no completed frozen repeated comparison. Earlier Paris itinerary
        and labour-law scheduling examples remain prospective. The preprocessing toolkit reproduces official
        totals on real public data, which is narrower evidence than autonomous acquisition and modelling. In one run on real Taipei data, the local Data Adapter produced a valid instance but
        ignored the toolkit, wrote its own transformations and added licence information that was not in the
        sources; the provenance record must therefore be written by code, not by the model.
  \item \textbf{Checker independence is limited.} Separate prompts and computations reduce direct reuse of
        Adapter or model code, but the roles share a base LLM and schema and can share the same error.
        Counts and totals do not establish per-entity data correctness, and sampled requirement checks do not
        verify every feasible state. The external judge remains case-specific and can itself fail.
  \item \textbf{Agent and retrieval effects remain hypotheses.} The direct baseline differs in tools,
        acquisition limits and feedback, so it does not isolate agent count. The new modelling-pattern
        retrieval ablation is implemented but has no completed repeated result. Earlier null results for
        splitting tuning roles or retrieving tuning knowledge do not answer the modelling-stage question.
  \item \textbf{Hardware evidence remains small.} Subsequent CGA and Leap Hybrid BQM comparisons use three repetitions per
        condition. Hybrid service results are distinct from direct QPU sampling, and use a ten-second limit
        versus the three-second CGA ceiling; the resource allocations do not establish hardware superiority.
        Effects depend on the problem; Fujitsu DA rejected a submission due to its monthly
        account limit. The earlier simulation advantage is small, and neither role splitting nor retrieval
        helped there. Large-scale repeated hardware and held-out backend-transfer studies remain open.
  \item \textbf{Functional coverage is not end-to-end reliability.} The integration checks in
        \cref{sec:integrationchecks} exercise parameterized compilation, local search and recovery fixtures.
        They do not establish success rates for live data acquisition, LLM diagnosis or model repair, and the
        new penalty policies have an eight-instance repeated simulation comparison, in which defaults had the
        lowest mean gap, plus small actual-CGA comparisons. A feasible
        baseline gives a comparison point, not a guarantee of application-level quality.
  \item \textbf{Recovery has a limited scope.} Checkpoint resume supports the four history-based penalty
        policies on typed models, not the existing agent or joint Optuna state. Supported searches can resume
        through the UI or command-line tool. An interrupted external solver call may
        require inspection before resubmission. Requirement coverage and formula-change checks cannot prove
        correctness of the original model, and a single-flip penalty floor cannot certify every global
        optimum. Data substitutions and requirement changes still need user judgment.
  \item \textbf{Small local model.} The local model fails as a modeller in recognisable ways (\cref{sec:mechanisms});
        our checks catch structural failures but not every semantic one: a model may pass all checks and still
        replace the requested objective with a weaker surrogate, which only the user's review catches. Data
        interpretation is also fallible (a labour-law rule of two rest days per seven was simplified to one).
  \item \textbf{Scale.} Brute-force verification is limited to about 26 variables; larger encodings are reported
        as unverifiable. Routing on the GPU annealer degrades beyond about 30 stops; multi-vehicle routing
        produced the weakest models of all categories, and connectivity constraints (network design) remain hard
        to express.
  \item \textbf{Benchmark scope.} The 100 benchmarks have tiny instances (median five variables), and our
        12-category suite uses synthetic data. Neither measures real-world difficulty.
  \item \textbf{Non-determinism.} Even at temperature zero, 27\% of benchmark outcomes changed between identical
        runs; all comparisons in this paper therefore use repeated runs for local conditions.
\end{itemize}
